% assets/w04-antiwindup-two.svg — 되감기 두 방식은 적분기 입력에서만 갈라진다.
%
%   bash _tools/tikz2svg.sh assets/src/w04-antiwindup-two.tex assets/w04-antiwindup-two.svg
%
% 4주차 1-12 의 그림. 대학원 볼트 figures/src/w03-antiwindup-two.tex 를 캡스톤 기호·
% 한국어로 다시 쓴 것.
%
% 두 방식은 루프의 나머지가 전부 같다. 다른 것은 적분기 바로 앞에서 무엇이 일어나는가
% 하나뿐이므로 그 자리만 나란히 놓고 나머지는 같은 모양으로 그린다. 포화 뒤가 왜
% 끊기는가(전체 루프)는 w04-windup.svg 가 이미 그렸다 — 여기서 플랜트와 바깥
% 되먹임을 다시 그리지 않는 이유다.
%
% 부호는 1-12 결과식과 같다 —  Idot = K_i e + K_b (tau_sat - tau_raw).
% 따라서 eps = tau_sat - tau_raw 이고 상한 포화 중에는 eps < 0 이다.
% 본 과목 모델 PID_byhand 는 (b) 되감기(back-calculation) 를 쓴다.
%
% 폭 약 176 mm. 본문 폭 186 mm 안.
\documentclass[border=6pt]{standalone}
% capstone-style.tex — 이 볼트의 TikZ 그림이 공유하는 스타일
%
% 쓰는 법:  % capstone-style.tex — 이 볼트의 TikZ 그림이 공유하는 스타일
%
% 쓰는 법:  % capstone-style.tex — 이 볼트의 TikZ 그림이 공유하는 스타일
%
% 쓰는 법:  \input{capstone-style}   (assets/src/ 안에서)
% 빌드:     bash _tools/tikz2svg.sh assets/src/w02-node-topic.tex assets/w02-node-topic.svg
%
% 왜 TikZ 인가
%   손으로 SVG 좌표를 쓰면 화살촉·선 굵기·글자 크기가 그림마다 미묘하게 달라진다.
%   그것이 "자료가 어설퍼 보인다"의 정체다. 스타일을 한 번만 정하고 상대좌표로
%   배치하면 좌표를 고쳐도 그림이 무너지지 않는다.
%   수식은 본문과 같은 TeX 글꼴로 나온다.
%
% 한글
%   ko.TeX(DVI 경로)를 쓴다. XeLaTeX 은 TikZ 그래픽을 PDF special 로 내보내는데
%   dvisvgm 이 그것을 받으려면 Ghostscript < 10.01 이 필요하다. 이 PC 에는
%   10.07 이 깔려 있어 그 경로가 막힌다. latex -> DVI -> dvisvgm 이 유일하게
%   동작하는 조합이다.

\usepackage[hangul]{kotex}
\usepackage{tikz}
\usepackage{amsmath}
\usepackage{xcolor}
\usetikzlibrary{arrows.meta, positioning, calc, fit, backgrounds, decorations.pathmorphing}

% ---- 색 --------------------------------------------------------------------
% 강조는 하나만. 나머지는 검정. 흑백 인쇄에서도 읽혀야 한다.
\definecolor{ink}{HTML}{1A1A1A}
\definecolor{soft}{HTML}{666666}
\definecolor{accent}{HTML}{7C3AED}
\definecolor{warn}{HTML}{B91C1C}

% 계층 색 — Simulink 모델의 블록 색과 같은 규칙을 쓴다
\definecolor{guid}{HTML}{1D4ED8}    % 유도
\definecolor{ctrl}{HTML}{EA580C}    % 제어
\definecolor{thr}{HTML}{D97706}     % 추진기
\definecolor{plant}{HTML}{16A34A}   % 운동모델
\definecolor{nav}{HTML}{7C3AED}     % 항법 · 신호처리
\definecolor{pathc}{HTML}{0D9488}   % 계획 경로 · 항적
\definecolor{errc}{HTML}{C2410C}    % 오차
\definecolor{flow}{HTML}{0369A1}    % 조류 · 바람

% ---- 선과 화살촉 -----------------------------------------------------------
\tikzset{
  sig/.style   = {ink, line width=0.5pt, -{Stealth[length=4pt, width=3pt]}},
  wire/.style  = {ink, line width=0.5pt},
  fb/.style    = {ink, line width=0.5pt, densely dashed,
                  -{Stealth[length=4pt, width=3pt]}},
  asig/.style  = {accent, line width=0.5pt, densely dashed,
                  -{Stealth[length=4pt, width=3pt]}},
  % 블록 — 흰 바탕, 직각, 검은 테두리. 색은 테두리에만 준다
  blk/.style   = {draw=ink, line width=0.55pt, fill=white,
                  minimum width=18mm, minimum height=8.5mm, inner sep=2pt, font=\small},
  gblk/.style  = {blk, draw=guid},
  cblk/.style  = {blk, draw=ctrl},
  tblk/.style  = {blk, draw=thr},
  pblk/.style  = {blk, draw=plant},
  nblk/.style  = {blk, draw=nav},
  % 합산점 · 분기점
  sum/.style   = {draw=ink, line width=0.55pt, fill=white, circle,
                  inner sep=0pt, minimum size=4.6mm, font=\scriptsize},
  dot/.style   = {ink, fill=ink, circle, inner sep=0pt, minimum size=1.5mm},
  % 글자
  sname/.style = {font=\small, inner sep=1.5pt},
  note/.style  = {font=\scriptsize, soft, inner sep=1.5pt},
  wnote/.style = {font=\scriptsize, warn, inner sep=1.5pt},
  anote/.style = {font=\scriptsize, accent, inner sep=1.5pt},
  axis/.style  = {ink, line width=0.5pt},
  tick/.style  = {font=\scriptsize, soft},
  % 묶음 상자 — 서브시스템 표시
  grp/.style   = {draw=soft, line width=0.4pt, densely dotted, rounded corners=1mm,
                  inner sep=3mm},
}
   (assets/src/ 안에서)
% 빌드:     bash _tools/tikz2svg.sh assets/src/w02-node-topic.tex assets/w02-node-topic.svg
%
% 왜 TikZ 인가
%   손으로 SVG 좌표를 쓰면 화살촉·선 굵기·글자 크기가 그림마다 미묘하게 달라진다.
%   그것이 "자료가 어설퍼 보인다"의 정체다. 스타일을 한 번만 정하고 상대좌표로
%   배치하면 좌표를 고쳐도 그림이 무너지지 않는다.
%   수식은 본문과 같은 TeX 글꼴로 나온다.
%
% 한글
%   ko.TeX(DVI 경로)를 쓴다. XeLaTeX 은 TikZ 그래픽을 PDF special 로 내보내는데
%   dvisvgm 이 그것을 받으려면 Ghostscript < 10.01 이 필요하다. 이 PC 에는
%   10.07 이 깔려 있어 그 경로가 막힌다. latex -> DVI -> dvisvgm 이 유일하게
%   동작하는 조합이다.

\usepackage[hangul]{kotex}
\usepackage{tikz}
\usepackage{amsmath}
\usepackage{xcolor}
\usetikzlibrary{arrows.meta, positioning, calc, fit, backgrounds, decorations.pathmorphing}

% ---- 색 --------------------------------------------------------------------
% 강조는 하나만. 나머지는 검정. 흑백 인쇄에서도 읽혀야 한다.
\definecolor{ink}{HTML}{1A1A1A}
\definecolor{soft}{HTML}{666666}
\definecolor{accent}{HTML}{7C3AED}
\definecolor{warn}{HTML}{B91C1C}

% 계층 색 — Simulink 모델의 블록 색과 같은 규칙을 쓴다
\definecolor{guid}{HTML}{1D4ED8}    % 유도
\definecolor{ctrl}{HTML}{EA580C}    % 제어
\definecolor{thr}{HTML}{D97706}     % 추진기
\definecolor{plant}{HTML}{16A34A}   % 운동모델
\definecolor{nav}{HTML}{7C3AED}     % 항법 · 신호처리
\definecolor{pathc}{HTML}{0D9488}   % 계획 경로 · 항적
\definecolor{errc}{HTML}{C2410C}    % 오차
\definecolor{flow}{HTML}{0369A1}    % 조류 · 바람

% ---- 선과 화살촉 -----------------------------------------------------------
\tikzset{
  sig/.style   = {ink, line width=0.5pt, -{Stealth[length=4pt, width=3pt]}},
  wire/.style  = {ink, line width=0.5pt},
  fb/.style    = {ink, line width=0.5pt, densely dashed,
                  -{Stealth[length=4pt, width=3pt]}},
  asig/.style  = {accent, line width=0.5pt, densely dashed,
                  -{Stealth[length=4pt, width=3pt]}},
  % 블록 — 흰 바탕, 직각, 검은 테두리. 색은 테두리에만 준다
  blk/.style   = {draw=ink, line width=0.55pt, fill=white,
                  minimum width=18mm, minimum height=8.5mm, inner sep=2pt, font=\small},
  gblk/.style  = {blk, draw=guid},
  cblk/.style  = {blk, draw=ctrl},
  tblk/.style  = {blk, draw=thr},
  pblk/.style  = {blk, draw=plant},
  nblk/.style  = {blk, draw=nav},
  % 합산점 · 분기점
  sum/.style   = {draw=ink, line width=0.55pt, fill=white, circle,
                  inner sep=0pt, minimum size=4.6mm, font=\scriptsize},
  dot/.style   = {ink, fill=ink, circle, inner sep=0pt, minimum size=1.5mm},
  % 글자
  sname/.style = {font=\small, inner sep=1.5pt},
  note/.style  = {font=\scriptsize, soft, inner sep=1.5pt},
  wnote/.style = {font=\scriptsize, warn, inner sep=1.5pt},
  anote/.style = {font=\scriptsize, accent, inner sep=1.5pt},
  axis/.style  = {ink, line width=0.5pt},
  tick/.style  = {font=\scriptsize, soft},
  % 묶음 상자 — 서브시스템 표시
  grp/.style   = {draw=soft, line width=0.4pt, densely dotted, rounded corners=1mm,
                  inner sep=3mm},
}
   (assets/src/ 안에서)
% 빌드:     bash _tools/tikz2svg.sh assets/src/w02-node-topic.tex assets/w02-node-topic.svg
%
% 왜 TikZ 인가
%   손으로 SVG 좌표를 쓰면 화살촉·선 굵기·글자 크기가 그림마다 미묘하게 달라진다.
%   그것이 "자료가 어설퍼 보인다"의 정체다. 스타일을 한 번만 정하고 상대좌표로
%   배치하면 좌표를 고쳐도 그림이 무너지지 않는다.
%   수식은 본문과 같은 TeX 글꼴로 나온다.
%
% 한글
%   ko.TeX(DVI 경로)를 쓴다. XeLaTeX 은 TikZ 그래픽을 PDF special 로 내보내는데
%   dvisvgm 이 그것을 받으려면 Ghostscript < 10.01 이 필요하다. 이 PC 에는
%   10.07 이 깔려 있어 그 경로가 막힌다. latex -> DVI -> dvisvgm 이 유일하게
%   동작하는 조합이다.

\usepackage[hangul]{kotex}
\usepackage{tikz}
\usepackage{amsmath}
\usepackage{xcolor}
\usetikzlibrary{arrows.meta, positioning, calc, fit, backgrounds, decorations.pathmorphing}

% ---- 색 --------------------------------------------------------------------
% 강조는 하나만. 나머지는 검정. 흑백 인쇄에서도 읽혀야 한다.
\definecolor{ink}{HTML}{1A1A1A}
\definecolor{soft}{HTML}{666666}
\definecolor{accent}{HTML}{7C3AED}
\definecolor{warn}{HTML}{B91C1C}

% 계층 색 — Simulink 모델의 블록 색과 같은 규칙을 쓴다
\definecolor{guid}{HTML}{1D4ED8}    % 유도
\definecolor{ctrl}{HTML}{EA580C}    % 제어
\definecolor{thr}{HTML}{D97706}     % 추진기
\definecolor{plant}{HTML}{16A34A}   % 운동모델
\definecolor{nav}{HTML}{7C3AED}     % 항법 · 신호처리
\definecolor{pathc}{HTML}{0D9488}   % 계획 경로 · 항적
\definecolor{errc}{HTML}{C2410C}    % 오차
\definecolor{flow}{HTML}{0369A1}    % 조류 · 바람

% ---- 선과 화살촉 -----------------------------------------------------------
\tikzset{
  sig/.style   = {ink, line width=0.5pt, -{Stealth[length=4pt, width=3pt]}},
  wire/.style  = {ink, line width=0.5pt},
  fb/.style    = {ink, line width=0.5pt, densely dashed,
                  -{Stealth[length=4pt, width=3pt]}},
  asig/.style  = {accent, line width=0.5pt, densely dashed,
                  -{Stealth[length=4pt, width=3pt]}},
  % 블록 — 흰 바탕, 직각, 검은 테두리. 색은 테두리에만 준다
  blk/.style   = {draw=ink, line width=0.55pt, fill=white,
                  minimum width=18mm, minimum height=8.5mm, inner sep=2pt, font=\small},
  gblk/.style  = {blk, draw=guid},
  cblk/.style  = {blk, draw=ctrl},
  tblk/.style  = {blk, draw=thr},
  pblk/.style  = {blk, draw=plant},
  nblk/.style  = {blk, draw=nav},
  % 합산점 · 분기점
  sum/.style   = {draw=ink, line width=0.55pt, fill=white, circle,
                  inner sep=0pt, minimum size=4.6mm, font=\scriptsize},
  dot/.style   = {ink, fill=ink, circle, inner sep=0pt, minimum size=1.5mm},
  % 글자
  sname/.style = {font=\small, inner sep=1.5pt},
  note/.style  = {font=\scriptsize, soft, inner sep=1.5pt},
  wnote/.style = {font=\scriptsize, warn, inner sep=1.5pt},
  anote/.style = {font=\scriptsize, accent, inner sep=1.5pt},
  axis/.style  = {ink, line width=0.5pt},
  tick/.style  = {font=\scriptsize, soft},
  % 묶음 상자 — 서브시스템 표시
  grp/.style   = {draw=soft, line width=0.4pt, densely dotted, rounded corners=1mm,
                  inner sep=3mm},
}

\tikzset{
  % 두 판을 나란히 놓아야 하므로 기본 18mm 블록은 넓다
  sblk/.style  = {draw=ink, line width=0.55pt, fill=white,
                  minimum width=11mm, minimum height=7.5mm, inner sep=1.5pt, font=\small},
  ablk/.style  = {draw=accent, line width=0.7pt, fill=white,
                  minimum width=13mm, minimum height=8.5mm, inner sep=1.5pt,
                  font=\small, text=accent},
  asum/.style  = {draw=accent, line width=0.55pt, fill=white, circle,
                  inner sep=0pt, minimum size=4.6mm, font=\scriptsize},
  awire/.style = {accent, line width=0.5pt},
}

\begin{document}
\begin{tikzpicture}[x=1mm, y=1mm]

\node[font=\small\bfseries, anchor=north west] at (-12,36)
     {되감기 두 방식 — 갈라지는 자리는 적분기 입력 한 곳뿐이다};
\node[anchor=north west, align=left, text width=176mm, font=\scriptsize, color=soft] at (-12,31) {%
  보라색 블록 바깥은 두 방식이 완전히 같다. 두 방식이 재는 넘친 양
  $\varepsilon = \tau_{\text{sat}} - \tau_{\text{raw}}$ 까지 같다 (상한 포화 중에는 음수).
  갈라지는 것은 그것을 적분기 입력에서 어떻게 쓰는가 하나뿐이다.};

% =======================================================================
% (a) 잠그기 (clamping)
% =======================================================================
\begin{scope}[shift={(0,0)}]
  \node[font=\small\bfseries, anchor=west] at (-10,15) {(a) \ \ 잠그기 (clamping)};
  \node[note, anchor=west, text width=78mm] at (-10,11)
       {더 쌓이기만 할 상황이면 적분기를 아예 \emph{끈다}};

  \node[sblk] (kp)  at (8,4)    {$K_p$};
  \node[sblk] (ki)  at (8,-8)   {$K_i$};
  \node[ablk] (gate) at (26,-8) {};
  \node[sblk] (int) at (44,-8)  {$1/s$};
  \node[sum]  (sb)  at (58,0)   {};
  \node[sblk] (sat) at (70,0)   {};

  % 스위치 글리프 — 열린 접점
  \draw[accent, line width=0.6pt] ([shift={(-4mm,0)}]gate.center) -- ([shift={(-1.6mm,0)}]gate.center);
  \draw[accent, line width=0.6pt] ([shift={( 1.6mm,0)}]gate.center) -- ([shift={( 4mm,0)}]gate.center);
  \draw[accent, line width=0.7pt] ([shift={(-1.6mm,0)}]gate.center) -- ([shift={( 1.3mm,2.3mm)}]gate.center);
  \fill[accent] ([shift={(-1.6mm,0)}]gate.center) circle (0.42);
  \fill[accent] ([shift={( 1.6mm,0)}]gate.center) circle (0.42);

  % 포화 글리프
  \draw[wire] ([shift={(-4.5mm,-2.1mm)}]sat.center) -- ([shift={(-1.7mm,-2.1mm)}]sat.center)
              -- ([shift={( 1.7mm, 2.1mm)}]sat.center) -- ([shift={( 4.5mm, 2.1mm)}]sat.center);

  % 배선
  \coordinate (br) at (-2,-2);
  \draw[sig] (-10,-2) -- node[sname, above, pos=0.15] {$e$} (br);
  \node[dot] at (br) {};
  \draw[wire] (br) |- (kp.west);   \draw[wire] (br) |- (ki.west);
  \draw[sig]  (ki) -- (gate);
  \draw[sig]  (gate) -- (int);
  \draw[wire] (kp) -| (sb);
  \draw[wire] (int) -| (sb);
  \draw[sig]  (sb) -- (sat);
  \draw[sig]  (sat) -- (81,0);
  \node[sname, anchor=west] at (81,0) {$\tau_{\text{sat}}$};
  \node[sname, anchor=east] at (65.2,-11) {$\tau_{\text{raw}}$};

  % eps 를 재는 자리
  \coordinate (ta) at (66,0);   \coordinate (tb) at (78,0);
  \node[dot] at (ta) {};  \node[dot] at (tb) {};
  \node[asum] (ep) at (72,-20) {};
  \node[anote, anchor=north] at (72,-23.5) {$\varepsilon$};
  \draw[awire] (ta) |- (ep);
  \draw[awire] (tb) |- (ep);
  \node[anote, anchor=east] at (69.4,-17.4) {$-$};
  \node[anote, anchor=west] at (74.6,-17.4) {$+$};

  % 게이트로 가는 논리
  \draw[asig] (ep) -| (26,-14.5);
  \node[anote, anchor=north, align=center, text width=60mm] at (30,-27)
       {$\varepsilon \neq 0$ 이고 오차가 같은 쪽을 더 밀고 있으면 연다 \\
        $\Rightarrow\ \dot I = 0$ — 기억을 얼린다};
\end{scope}

% =======================================================================
% (b) 되감기 (back-calculation) — 본 과목이 쓰는 방식
% =======================================================================
\begin{scope}[shift={(104,0)}]
  \node[font=\small\bfseries, anchor=west] at (-10,15) {(b) \ \ 되감기 (back-calculation)};
  \node[note, anchor=west, text width=78mm] at (-10,11)
       {넘친 양을 \emph{되돌려} 적분기를 끌어낸다 — 본 과목 \textsf{PID\_byhand} 가 이것};

  \node[sblk] (kp)  at (8,4)    {$K_p$};
  \node[sblk] (ki)  at (8,-8)   {$K_i$};
  \node[asum] (bs)  at (26,-8)  {};
  \node[sblk] (int) at (44,-8)  {$1/s$};
  \node[sum]  (sb)  at (58,0)   {};
  \node[sblk] (sat) at (70,0)   {};
  \node[ablk] (kaw) at (26,-20) {$K_b$};

  \draw[wire] ([shift={(-4.5mm,-2.1mm)}]sat.center) -- ([shift={(-1.7mm,-2.1mm)}]sat.center)
              -- ([shift={( 1.7mm, 2.1mm)}]sat.center) -- ([shift={( 4.5mm, 2.1mm)}]sat.center);

  \coordinate (br) at (-2,-2);
  \draw[sig] (-10,-2) -- node[sname, above, pos=0.15] {$e$} (br);
  \node[dot] at (br) {};
  \draw[wire] (br) |- (kp.west);   \draw[wire] (br) |- (ki.west);
  \draw[sig]  (ki) -- (bs);
  \draw[sig]  (bs) -- (int);
  \draw[wire] (kp) -| (sb);
  \draw[wire] (int) -| (sb);
  \draw[sig]  (sb) -- (sat);
  \draw[sig]  (sat) -- (81,0);
  \node[sname, anchor=west] at (81,0) {$\tau_{\text{sat}}$};
  \node[sname, anchor=east] at (65.2,-11) {$\tau_{\text{raw}}$};

  \coordinate (ta) at (66,0);   \coordinate (tb) at (78,0);
  \node[dot] at (ta) {};  \node[dot] at (tb) {};
  \node[asum] (ep) at (72,-20) {};
  \node[anote, anchor=north] at (72,-23.5) {$\varepsilon$};
  \draw[awire] (ta) |- (ep);
  \draw[awire] (tb) |- (ep);
  \node[anote, anchor=east] at (69.4,-17.4) {$-$};
  \node[anote, anchor=west] at (74.6,-17.4) {$+$};

  % eps -> Kb -> 적분기 입력에 더한다
  \draw[asig] (ep) -- (kaw);
  \draw[asig] (kaw) -- (bs);
  \node[anote, anchor=east] at (24.4,-11.5) {$+$};
  \node[anote, anchor=north, align=center, text width=60mm] at (30,-27)
       {$\dot I = K_i e + K_b\,\varepsilon$ — 끊임없이 끌어당긴다 \\
        멈추는 자리는 0 이 아니라 $\varepsilon = -(K_i/K_b)\,e$};
\end{scope}

\end{tikzpicture}
\end{document}
